% This document was generated by an LLM.

\documentclass{essay}

\title{Self-Check Questions:\\ Power Series and Coefficient Matching}
\author{}
\date{}

\begin{document}

\maketitle

\begin{enumerate}[leftmargin=*, itemsep=8pt]

  \item What does it mean for two power series to be equal?
    \begin{answer}
      Two power series $\sum_{n=0}^{\infty} a_n x^n$ and
      $\sum_{n=0}^{\infty} b_n x^n$ are equal if they converge to the
      same value for every $x$ in some open interval $(-R, R)$ with
      $R > 0$. By the uniqueness theorem this forces $a_n = b_n$ for
      all $n \ge 0$.
    \end{answer}

  \item Under what conditions can two power series be compared
    coefficient-by-coefficient?
    \begin{answer}
      Both series must (i) converge on a common open interval
      containing $0$ and (ii) represent the same function on that
      interval. Under these conditions the uniqueness theorem
      guarantees that corresponding coefficients are equal.
    \end{answer}

  \item Why is the statement \[ \sum_{n=0}^{\infty} c_n x^n = 0 \]
    stronger than merely saying that the series equals zero at \(x=0\)?
    \begin{answer}
      The equation at $x = 0$ alone yields only $c_0 = 0$. The
      identity as a power series asserts that the sum is zero
      \emph{for every} $x$ in some interval around $0$, which (by the
      uniqueness theorem) forces every coefficient $c_n$ to vanish.
    \end{answer}

  \item State the uniqueness theorem for power series.
    \begin{answer}
      If $\displaystyle\sum_{n=0}^{\infty} a_n x^n =
      \sum_{n=0}^{\infty} b_n x^n$ for all $x$ in some open interval
      containing $0$, then $a_n = b_n$ for all $n \ge 0$.
    \end{answer}

  \item How can the uniqueness theorem be proved using derivatives?
    \begin{answer}
      Let $f(x) = \sum_{n=0}^{\infty} c_n x^n = 0$ on $(-R,R)$.
      Setting $x = 0$ gives $c_0 = 0$. A convergent power series may
      be differentiated term-by-term inside its interval of
      convergence, so $f'(x) = \sum_{n=1}^{\infty} n c_n x^{n-1} =
      0$. Setting $x = 0$ gives $c_1 = 0$. Repeating, the $k$-th
      derivative satisfies $f^{(k)}(0) = k!\, c_k = 0$, hence $c_k =
      0$ for every $k \ge 0$.
    \end{answer}

  \item Why does \[ \sum_{n=0}^{\infty} c_n x^n = 0 \] imply \(c_0=0\)?
    \begin{answer}
      Substitute $x = 0$. Every term with $n \ge 1$ contains a factor
      of $x$ and therefore vanishes, leaving $c_0 \cdot 1 = 0$, so $c_0 = 0$.
    \end{answer}

  \item Why does differentiating the series imply \(c_1=0\)?
    \begin{answer}
      Differentiating term-by-term gives $\sum_{n=1}^{\infty} n c_n
      x^{n-1} = 0$. Setting $x = 0$ in this new series yields $1
      \cdot c_1 = 0$, so $c_1 = 0$.
    \end{answer}

  \item What is the general formula for \(c_n\) in terms of
    derivatives of the sum at \(x=0\)?
    \begin{answer}
      If $f(x) = \sum_{n=0}^{\infty} c_n x^n$, then differentiating
      $n$ times and evaluating at $x = 0$ gives $f^{(n)}(0) = n!\,
      c_n$, so \[ c_n = \frac{f^{(n)}(0)}{n!}. \]
    \end{answer}

  \item Rewrite \[ \sum_{n=1}^{\infty} a_n x^{n-1} \] as a power
    series in \(x^n\).
    \begin{answer}
      Substitute $m = n - 1$ (so $n = m + 1$); when $n = 1$, $m = 0$:
      \[
        \sum_{n=1}^{\infty} a_n x^{n-1}
        = \sum_{m=0}^{\infty} a_{m+1}\, x^m.
      \]
      Relabelling $m$ back to $n$: $\displaystyle\sum_{n=0}^{\infty}
      a_{n+1}\, x^n$.
    \end{answer}

  \item Why is reindexing necessary before matching coefficients?
    \begin{answer}
      Coefficient matching works by equating the multipliers of the
      \emph{same} power $x^n$ on both sides. If the series are
      expressed in different powers (e.g.\ $x^{n-1}$ versus $x^n$)
      the $k$-th summand of each series corresponds to a different
      power of $x$, so naive comparison produces a wrong (or
      meaningless) relation.
    \end{answer}

  \item What mistake occurs if two sums involve different powers of
    \(x\) and you compare coefficients immediately?
    \begin{answer}
      You would equate the coefficient of $x^n$ in one series against
      the coefficient of $x^n$ in the other, but one series might
      have its $n$-th term equal to $x^{n-1}$ rather than $x^n$. The
      resulting ``recurrence'' would relate coefficients that
      actually belong to different powers of $x$, giving an incorrect equation.
    \end{answer}

  \item Given \[ \sum_{n=0}^{\infty} a_n x^n + \sum_{n=2}^{\infty}
    b_n x^n, \] what is the coefficient of \(x^1\)?
    \begin{answer}
      The first series contributes $a_1$ for the power $x^1$. The
      second series starts at $n = 2$, so it has no $x^1$ term. The
      coefficient of $x^1$ is $a_1$.
    \end{answer}

  \item What is the coefficient of \(x^5\)?
    \begin{answer}
      Both series include $n = 5$, so the coefficient of $x^5$ is $a_5 + b_5$.
    \end{answer}

  \item Suppose \[ \sum_{n=0}^{\infty} a_n x^n + \sum_{n=3}^{\infty}
    b_n x^n =0. \] Why must \(a_0,a_1,a_2\) be treated separately?
    \begin{answer}
      The second series has no terms for $n = 0, 1, 2$. For those
      powers of $x$,
      the only contribution to the combined sum comes from the first
      series alone, so
      the equation $(\cdots) = 0$ for each of $x^0$, $x^1$, $x^2$ involves only
      $a_0$, $a_1$, $a_2$ respectively, and each must vanish on its own.
    \end{answer}

  \item What equation do you obtain from the coefficient of \(x^0\)?
    \begin{answer}
      $a_0 = 0$.
    \end{answer}

  \item What equation do you obtain from the coefficient of \(x^1\)?
    \begin{answer}
      $a_1 = 0$.
    \end{answer}

  \item What equation do you obtain from the coefficient of \(x^2\)?
    \begin{answer}
      $a_2 = 0$.
    \end{answer}

  \item What equation do you obtain for \(n \ge 3\)?
    \begin{answer}
      Both series contribute, so $a_n + b_n = 0$, i.e.\ $b_n = -a_n$ for all
      $n \ge 3$.
    \end{answer}

  \item How would the answer change if the second series started at \(n=5\)?
    \begin{answer}
      Then $a_0, a_1, a_2, a_3, a_4$ would each be forced to zero individually
      (isolated terms), and for $n \ge 5$ one would obtain $a_n + b_n = 0$.
    \end{answer}

  \item If \[ y=\sum_{n=0}^{\infty} a_n x^n, \] what is \(y'\)?
    \begin{answer}
      \[
        y' = \sum_{n=1}^{\infty} n\, a_n\, x^{n-1}.
      \]
    \end{answer}

  \item What is \(y''\)?
    \begin{answer}
      \[
        y'' = \sum_{n=2}^{\infty} n(n-1)\, a_n\, x^{n-2}.
      \]
    \end{answer}

  \item Why do the sums for \(y\), \(y'\), and \(y''\) begin at
    different indices?
    \begin{answer}
      Differentiating a constant term produces zero, so it drops out of the sum.
      The constant term ($n=0$) of $y$ vanishes in $y'$, forcing $y'$
      to start at
      $n=1$. The $n=0$ and $n=1$ terms of $y'$ collectively
      contribute at most a
      constant, which vanishes in $y''$, so $y''$ starts at $n=2$.
    \end{answer}

  \item After differentiating, how do you reindex so that all series
    are expressed in powers of \(x^n\)?
    \begin{answer}
      For $y'$, substitute $m = n-1$ ($n = m+1$, lower limit $m=0$):
      \[
        y' = \sum_{m=0}^{\infty} (m+1)\,a_{m+1}\,x^m.
      \]
      For $y''$, substitute $m = n-2$ ($n = m+2$, lower limit $m=0$):
      \[
        y'' = \sum_{m=0}^{\infty} (m+2)(m+1)\,a_{m+2}\,x^m.
      \]
      Rename $m \to n$ in both to obtain all three series in powers
      of $x^n$ starting
      at $n=0$.
    \end{answer}

  \item Why is matching coefficients impossible until all sums use
    the same power \(x^n\)?
    \begin{answer}
      The uniqueness theorem equates multipliers of identical powers
      of $x$. If one
      series runs in $x^n$ and another in $x^{n-2}$, the $k$-th term of each
      represents a different power of $x$, so there is no consistent meaning to
      ``the coefficient of $x^n$'' across the two sums until they are expressed
      uniformly.
    \end{answer}

  \item When solving an ODE by power series, what usually produces
    the recurrence relation?
    \begin{answer}
      After substituting the ansatz $y = \sum a_n x^n$, reindexing
      all series so they
      run in the same power $x^n$, and combining into a single sum
      equal to zero,
      equating the coefficient of $x^n$ to zero yields the recurrence
      relation, which
      expresses $a_{n+k}$ (for some fixed $k$) in terms of earlier coefficients.
    \end{answer}

  \item What usually determines the first few coefficients separately
    from the recurrence relation?
    \begin{answer}
      The initial conditions of the ODE --- typically $y(0)$ and $y'(0)$ for a
      second-order equation --- fix $a_0$ and $a_1$ directly. The recurrence
      relation then determines all subsequent coefficients in terms
      of these two.
      Additionally, isolated low-index terms (powers of $x$ not covered by the
      combined recurrence) must be set to zero individually,
      sometimes also pinning
      early coefficients.
    \end{answer}

  \item Why is coefficient matching essentially equivalent to
    evaluating derivatives at \(x=0\)?
    \begin{answer}
      The coefficient $c_n$ of a power series satisfies $c_n = f^{(n)}(0)/n!$.
      Requiring $\sum a_n x^n = \sum b_n x^n$ and then matching $a_n = b_n$ is
      therefore exactly the same as requiring $f^{(n)}(0) = g^{(n)}(0)$ for all
      $n \ge 0$, i.e.\ that all derivatives of the two functions agree at $0$.
    \end{answer}

  \item Could coefficient matching fail for a Taylor series outside
    its radius of convergence?
    \begin{answer} Yes. A Taylor series converges to the original
      function only inside the radius of convergence $R$. Outside
      $|x| < R$ the series may diverge or converge to a different
      value, so the identity theorem (which underpins coefficient
      matching) no longer applies, and matching coefficients yields
      no information about the function's behaviour there.
    \end{answer}

  \item Why is analyticity essential for the uniqueness theorem?
    \begin{answer}
      A smooth ($C^\infty$) function need not equal its Taylor series
      anywhere except at the centre. If $f$ is not analytic, two
      distinct functions can share every Taylor coefficient at $0$
      and still differ on every neighbourhood of $0$. Analyticity is
      precisely the condition that $f$ \emph{equals} its Taylor
      series in some disk, which is what makes the map $f \mapsto
      (c_n)$ injective and coefficient matching valid.
    \end{answer}

  \item Give an example of a smooth \((C^\infty)\) function whose
    Taylor series is identically zero but whose function is not.
    \begin{answer}
      Define
      \[
        f(x) =
        \begin{cases}
          e^{-1/x^2} & x \neq 0, \\
          0          & x = 0.
        \end{cases}
      \]
      Every derivative of $f$ at $x = 0$ equals $0$, so the Taylor
      series of $f$ centred at $0$ is $\sum 0 \cdot x^n \equiv 0$,
      yet $f(x) > 0$ for all $x \neq 0$.
    \end{answer}

  \item Why does such an example show that power series methods are
    fundamentally about analytic functions, not merely differentiable functions?
    \begin{answer}
      This function is $C^\infty$ but its Taylor series does not
      converge to it (the series is identically $0$ while the
      function is not). Any method that represents a solution as a
      power series and then extracts the solution by matching or
      summing those series is implicitly assuming the solution is
      analytic. For merely smooth functions the power series may be
      useless as a representation, and the method breaks down entirely.
    \end{answer}

  \item In a power series solution of an ODE, what role do the
    arbitrary constants \(a_0\) and \(a_1\) typically play for a
    second-order equation?
    \begin{answer}
      They correspond to the two constants of integration in the
      general solution. Concretely, $a_0 = y(0)$ and $a_1 = y'(0)$.
      Every higher coefficient $a_n$ ($n \ge 2$) is determined by the
      recurrence relation in terms of $a_0$ and $a_1$, so the general
      power series solution is a linear combination of two
      independent series, one proportional to $a_0$ and the other to $a_1$.
    \end{answer}

  \item Why do second-order ODEs usually leave two coefficients free?
    \begin{answer}
      By the existence-uniqueness theorem for linear ODEs, the
      solution space of a homogeneous second-order linear ODE is
      two-dimensional. In the power series approach this is reflected
      by the recurrence relation: it expresses $a_{n+2}$ in terms of
      $a_n$ and possibly $a_{n+1}$, advancing the index by $2$ (or by
      $1$), so the two ``seed'' values $a_0$ and $a_1$ propagate into
      the entire sequence and remain free parameters.
    \end{answer}

  \item How would you explain, in one paragraph, why ``equating
    coefficients'' is mathematically justified rather than a formal trick?
    \begin{answer}
      Equating coefficients is justified by the uniqueness theorem
      for analytic functions: if two power series converge to the
      same function on any open interval containing $0$, their
      coefficients must be identical. When we write an ODE as a
      single power series equal to zero, we are asserting that the
      resulting function is identically zero on some interval --- not
      just at one point. The uniqueness theorem then guarantees that
      each individual coefficient must be zero, so we are not merely
      ``pretending'' that powers of $x$ are independent; we are using
      the proved injectivity of the map from an analytic function to
      its Taylor-coefficient sequence. The method is therefore as
      rigorous as any other that invokes this theorem, provided the
      solution is known (or assumed) to be analytic in a
      neighbourhood of the expansion point.
    \end{answer}
\end{enumerate}

\end{document}
